% from ega4-2-fr.tex
\oldpage[IV-2]{74}

par $\mathscr J$ est alors géométriquement réduit et géométriquement
irréductible, donc géométriquement intègre. Inversement, si $Y$ est
géométriquement intègre, pour toute extension $k'$ de $k$,
$\operatorname{Supp}(\mathscr F')$ est irréductible, donc, si $X'$ est
localement noethérien, $\mathscr F'$ est intègre~; cela montre que c) entraîne
a), et que c) entraîne d) lorsque $X$ est localement de type fini sur $k$.
\end{proof}

\begin{definition}[4.6.19]
\label{IV.2.4.6.19-fr}
Lorsque $\mathscr F$ vérifie les conditions équivalentes de (4.6.18), on dit
que $\mathscr F$ est \emph{géométriquement intègre} sur $k$.
\end{definition}

\begin{proposition}[4.6.20]
\label{IV.2.4.6.20-fr}
Soient $k$ un corps, $X$ un $k$-préschéma localement noethérien, $\mathscr F$
un $\mathcal O_X$-Module cohérent. Soit $K$ une extension de $k$ telle que
$X\otimes_k K$ soit localement noethérien. Alors, si $\mathscr F$ est
géométriquement réduit (resp. géométriquement intègre) sur $k$,
$\mathscr F\otimes_k K$ est géométriquement réduit (resp. géométriquement
intègre) sur $K$.
\end{proposition}

\begin{proof}
En effet, pour toute extension finie $K'$ de $K$,
\[
(X\otimes_k K)\otimes_K K'=X\otimes_k K'
\]
est localement noethérien, et on a vu dans les démonstrations de (4.6.16) et
(4.6.18) que l'hypothèse sur $\mathscr F$ entraîne que
$\mathscr F\otimes_k K'$ est réduit (resp. intègre), d'où la conclusion.
\end{proof}

\begin{proposition}[4.6.21]
\label{IV.2.4.6.21-fr}
Soient $k$ un corps, $X$, $Y$ deux $k$-préschémas localement noethériens tels
que $X\times_k Y$ soit localement noethérien. Soient $\mathscr F$ un
$\mathcal O_X$-Module cohérent, $\mathscr G$ un $\mathcal O_Y$-Module cohérent.
\begin{enumerate}
\item[(i)] Si $\mathscr F$ est géométriquement réduit (resp. géométriquement
intègre) et $\mathscr G$ réduit (resp. intègre), alors
$\mathscr F\otimes_k\mathscr G$ est réduit (resp. intègre).
\item[(ii)] Si $\mathscr F$ et $\mathscr G$ sont géométriquement réduits (resp.
géométriquement intègres), il en est de même de
$\mathscr F\otimes_k\mathscr G$.
\end{enumerate}
\end{proposition}

\begin{proof}
\begin{enumerate}
\item[(i)] Utilisant (3.2.3), (4.6.16) et (4.6.18), on peut se borner au cas où
$\operatorname{Supp}(\mathscr F)=X$, $\operatorname{Supp}(\mathscr G)=Y$, $X$
étant géométriquement réduit (resp. géométriquement intègre), $Y$ réduit
(resp. intègre), $\mathscr F$ et $\mathscr G$ sans cycles premiers associés
immergés, et $\mathscr F_x$ étant isomorphe à $\mathcal O_x$ en tout point
maximal $x$ de $X$, $\mathscr G_y$ isomorphe à $\mathcal O_y$ en tout point
maximal $y$ de $Y$. On sait alors (4.2.3) que
$\mathscr F\otimes_k\mathscr G$ est sans cycle premier associé immergé, et que
les cycles premiers associés à $\mathscr F\otimes_k\mathscr G$ sont exactement
les composantes irréductibles de $X\times_k Y$ (4.2.5). On peut donc se borner
au cas où $X$ et $Y$ sont irréductibles~; alors, si $X$ est géométriquement
réduit et $Y$ réduit, on sait (4.2.4, (ii)) que $X\times_k Y$ est réduit~; si
$X$ est géométriquement intègre et $Y$ intègre, on sait que $X\times_k Y$ est
en outre irréductible (4.5.8, (i)), donc intègre. Enfin, tout point maximal $z$
de $X\times_k Y$ est au-dessus de $x$ et de $y$, donc les hypothèses sur
$\mathscr F$ et $\mathscr G$ entraînent, en vertu de (I, 9.1.12), que
$(\mathscr F\otimes_k\mathscr G)_z$ est isomorphe à $\mathcal O_z$, ce qui
achève la démonstration dans ce cas.

\item[(ii)] Le raisonnement est analogue, mais ici (après réduction au cas où
$X=\operatorname{Supp}(\mathscr F)$, $Y=\operatorname{Supp}(\mathscr G)$),
$X\times_k Y$ est géométriquement réduit (resp. géométriquement intègre) en
vertu de (4.6.5, (ii)).
\end{enumerate}
\end{proof}

Les définitions (4.6.18) et (4.6.19) se localisent~:

\begin{definition}[4.6.22]
\label{IV.2.4.6.22-fr}
Soient $k$ un corps, $X$ un $k$-préschéma localement noethérien, $\mathscr F$
un $\mathcal O_X$-Module cohérent. On dit que $\mathscr F$ est
\emph{géométriquement réduit} ou \emph{séparable} (resp. géométriquement
ponctuellement intègre) en un point $x\in X$ si, pour toute extension finie
radicielle (resp. finie) $k'$ de $k$, $\mathscr F\otimes_k k'$ est réduit
(resp. intègre) en chacun des points $x'$ de $X\otimes_k k'$ au-dessus de $x$
(cf. (3.2.2) et (3.2.4)).
\end{definition}

